% XeTeX/Tectonic has no TU/ptm definition; map IEEEtran's Times (ptm) to TeX Gyre Termes.
% Must precede \documentclass because IEEEtran selects its font while loading.
\ifdefined\XeTeXversion
  \DeclareFontFamily{TU}{ptm}{}
  \DeclareFontShape{TU}{ptm}{m}{n}{<->"[texgyretermes-regular.otf]:mapping=tex-text"}{}
  \DeclareFontShape{TU}{ptm}{m}{it}{<->"[texgyretermes-italic.otf]:mapping=tex-text"}{}
  \DeclareFontShape{TU}{ptm}{bx}{n}{<->"[texgyretermes-bold.otf]:mapping=tex-text"}{}
  \DeclareFontShape{TU}{ptm}{bx}{it}{<->"[texgyretermes-bolditalic.otf]:mapping=tex-text"}{}
  \DeclareFontShape{TU}{ptm}{b}{n}{<->ssub*ptm/bx/n}{}
  \DeclareFontShape{TU}{ptm}{b}{it}{<->ssub*ptm/bx/it}{}
  \DeclareFontShape{TU}{ptm}{m}{sl}{<->ssub*ptm/m/it}{}
  \DeclareFontShape{TU}{ptm}{bx}{sl}{<->ssub*ptm/bx/it}{}
  \DeclareFontShape{TU}{ptm}{m}{sc}{<->"[texgyretermes-regular.otf]:mapping=tex-text;+smcp"}{}
\fi
\documentclass[journal]{IEEEtran}

\usepackage{cite}
\usepackage{amsmath,amssymb,amsfonts}
\usepackage{graphicx}
\usepackage{textcomp}
\usepackage{xcolor}
\usepackage{booktabs}
\usepackage{multirow}
\ifdefined\XeTeXversion\else
  \usepackage[utf8]{inputenc}
  \usepackage[T1]{fontenc}
\fi
\usepackage{microtype}
\usepackage{array}
\usepackage{amsthm}
\usepackage{hyperref}

\newtheorem{assumption}{Assumption}
\newtheorem{definition}{Definition}
\newtheorem{theorem}{Theorem}
\newtheorem{proposition}{Proposition}
\newtheorem{corollary}{Corollary}
\newtheorem{lemma}{Lemma}
\newtheorem{remark}{Remark}

\interdisplaylinepenalty=2500
\emergencystretch=2em

\newcommand{\vx}{\mathbf{x}}
\newcommand{\vu}{\mathbf{u}}
\newcommand{\vs}{\mathbf{s}}
\newcommand{\vz}{\mathbf{z}}
\newcommand{\vc}{\mathbf{c}}
\newcommand{\vchi}{\boldsymbol{\chi}}
\newcommand{\ve}{\mathbf{e}}
\newcommand{\vr}{\mathbf{r}}
\newcommand{\vg}{\mathbf{g}}
\newcommand{\mA}{\mathbf{A}}
\newcommand{\mB}{\mathbf{B}}
\newcommand{\mC}{\mathbf{C}}
\newcommand{\mD}{\mathbf{D}}
\newcommand{\mF}{\mathbf{F}}
\newcommand{\mG}{\mathbf{G}}
\newcommand{\mQ}{\mathbf{Q}}
\newcommand{\mR}{\mathbf{R}}
\newcommand{\mP}{\mathbf{P}}
\newcommand{\mU}{\mathbf{U}}
\newcommand{\mV}{\mathbf{V}}
\newcommand{\mW}{\mathbf{W}}
\newcommand{\mX}{\mathbf{X}}
\newcommand{\mZ}{\mathbf{Z}}
\newcommand{\mI}{\mathbf{I}}
\newcommand{\mTheta}{\boldsymbol{\Theta}}
\newcommand{\R}{\mathbb{R}}
\newcommand{\norm}[1]{\left\lVert #1\right\rVert}
\newcommand{\argmin}{\mathop{\mathrm{arg\,min}}}
% Placeholder frame for a figure to be drawn later: \figph{width}{height}{description}
\newcommand{\figph}[3]{\fbox{\parbox[c][#2][c]{\dimexpr#1-2\fboxsep-2\fboxrule\relax}{\centering\small #3}}}

\begin{document}

\title{Transformer--Koopman: History-Conditioned Scheduled Prediction and Frozen-LTV Model Predictive Control}
\author{Bolin Tao}
\markboth{IEEE Journal Draft}{Tao: Transformer--Koopman}
\maketitle

\begin{abstract}
Koopman-inspired predictors provide a structured representation of nonlinear dynamics by evolving observables through finite-dimensional linear maps. A fixed lifted transition, however, can be restrictive when the measured state does not determine the relevant operating condition. This paper develops a history-conditioned scheduled predictor in which a temporal lifting map is complemented by a context-dependent low-rank transition and a measurement-dependent input modulation. The transition and input matrices are bounded functions of a history context and the current lifted measurement, respectively. For fixed scheduling signals the predictor is linear in the lifted state and input, while recursive scheduling remains generally nonlinear. Under common-history, work-domain, Lipschitz, and one-step residual conditions, we derive a finite-horizon open-loop prediction bound. For control, a nominal rollout supplies scheduling matrices that are frozen over the horizon, yielding an LTV prediction model and a convex quadratic program under polyhedral input constraints. We further derive a finite-horizon bound between the complete scheduled model and its frozen approximation, and state a conditional cost-suboptimality result relative to the global optimum of the complete problem.
\end{abstract}

\begin{IEEEkeywords}
Koopman modeling, temporal lifting, scheduled prediction, low-rank modulation, LTV-MPC, model predictive control.
\end{IEEEkeywords}

% =========================================================
\section{Introduction}
\label{sec:introduction}
% =========================================================

Predicting nonlinear dynamical systems from observed input--output trajectories is a central task in data-driven modeling. In many applications, the measured variables provide only an incomplete description of the physical state, and the response to an input depends on dynamical information accumulated over time. A predictor based solely on the current measurement may therefore be unable to distinguish trajectories that appear similar at one instant but evolve differently thereafter. At the same time, a highly expressive nonlinear model is not necessarily easy to propagate, analyze, or incorporate into subsequent decision-making. These considerations motivate predictive representations that exploit observation histories without abandoning a useful dynamical structure \cite{ljung1999system,nelles2020nonlinear,forgione2021dynonet}.

Koopman operator theory offers a framework for separating nonlinear state representation from the evolution of observables. Although the physical-state dynamics may be nonlinear, the Koopman operator acts linearly on observable functions \cite{koopman1931hamiltonian,mezic2005spectral,brunton2022modern}. Practical methods approximate this infinite-dimensional description using a finite set of observables, as in extended dynamic mode decomposition and controlled lifted predictors \cite{williams2015data,proctor2016dmdc,korda2018linear}. Neural Koopman models further learn the observable coordinates from data rather than prescribing them in advance \cite{takeishi2017koopman,lusch2018deep,otto2019linearly,han2020deep}. Their appeal lies in combining a flexible representation with a simple lifted evolution law. Nevertheless, finite-dimensional closure is generally approximate, and the quality of a learned predictor depends jointly on the information captured by its observables and the adequacy of its fitted evolution law.

These two requirements become particularly important when the measured coordinates do not determine the relevant dynamical condition. One approach enriches the representation: delay-coordinate constructions incorporate past observations into Koopman observables \cite{kamb2020timedelay,das2019delay}, while temporal neural encoders learn features from measurement sequences \cite{han2024robust,shenoy2025temporal}. Transformer-based Koopman forecasting models similarly exploit temporal dependencies within an observation window \cite{forootani2025deepkoopformer}. Such methods provide richer lifted coordinates, but a predictor using a single fixed matrix must still represent the behavior of different histories through the same finite-dimensional transition. A complementary approach adapts the transition itself. Koopa uses context-aware operators for nonstationary forecasting \cite{liu2023koopa}, and MamKO generates Koopman matrices from recent data for dynamical modeling and predictive control \cite{li2025mamko}. Thus, neither temporal lifting nor history-dependent operators are new in isolation. The relevant question is how their roles can be coordinated within a structured recursive predictor.

This question has consequences beyond the choice of network architecture. When the transition matrix is conditioned on a history window, a multi-step rollout must update that window using the model's own predicted measurements and the prescribed inputs. Errors then enter the prediction twice: directly through the propagated lifted state and indirectly through the context that schedules subsequent transitions. Consequently, an accurate one-step fit does not, by itself, characterize the behavior of the recursively scheduled model over a finite horizon. Moreover, unconstrained changes in learned transition matrices complicate the interpretation and analysis of the resulting predictor. A useful formulation should therefore make the roles of temporal representation and operator scheduling explicit, restrict the magnitude of scheduling-induced changes, and account for the feedback from prediction errors to future scheduling.

Motivated by these requirements, we develop a \emph{history-conditioned scheduled Koopman-inspired predictor}. A temporal Transformer extracts hidden lifted features from a finite window of measurements, while the latest measurement remains an explicit component of the lifted state. A separate context network processes recent measurements and past inputs to schedule the lifted transition matrix; the input matrix is modulated using the current measured component. Both matrices are expressed as nominal terms plus bounded low-rank corrections. This design distinguishes the information used to construct the lifted state from the information used to adapt its evolution, and it provides explicit bounds on the matrix modulation. The model is linear in the lifted state and input for prescribed scheduling signals, but its complete recursively scheduled evolution is generally nonlinear. We specify the associated history-update rule and derive a finite-horizon prediction-error bound that incorporates errors in the scheduling context under stated regularity, residual, and work-domain assumptions.

The same structure also permits a natural, but nontrivial, extension to predictive control. Because future scheduling signals depend on the predicted trajectory, directly optimizing inputs through the complete scheduled model generally produces a nonlinear program. We instead generate a nominal rollout, freeze the resulting scheduling matrices over the optimization horizon, and obtain a linear time-varying predictor. With the quadratic measurement-tracking objective and polyhedral input constraints considered here, the corresponding model predictive control (MPC) problem is a convex quadratic program. This freezing procedure is an approximation rather than an exact reformulation. We therefore bound the discrepancy between the recursively scheduled and frozen trajectories for a common candidate input sequence and obtain a conditional cost-suboptimality bound under a uniform approximation condition and a common feasible set. These results concern finite-horizon prediction and optimization approximation; they do not imply closed-loop stability or recursive feasibility.

The contributions of this paper are threefold:
\begin{enumerate}
    \item \emph{History-conditioned structured prediction:} We formulate a measurement-preserving temporal lifting together with a history-scheduled transition and a measurement-modulated input map. Bounded low-rank corrections define an explicit, analyzable parameterization of the lifted dynamics.
    \item \emph{Analysis of recursive history-dependent prediction:} We formalize the update of measurement and input histories during open-loop rollout and establish a finite-horizon error bound that accounts for the influence of prediction errors on subsequent operator scheduling.
    \item \emph{Frozen-LTV predictive control and approximation analysis:} We construct a nominal-rollout frozen-LTV MPC formulation that yields a convex quadratic program, quantify the finite-horizon error incurred by freezing the schedule, and establish a conditional control-cost bound relative to the complete scheduled optimization problem.
\end{enumerate}

The remainder of the paper is organized as follows. Section~\ref{sec:prelim} introduces the system setting and history-based prediction problem. Section~\ref{sec:model} develops the scheduled predictor and its error analysis. Section~\ref{sec:control} presents the frozen-LTV MPC formulation and approximation bounds. Section~\ref{sec:experiments} describes the experimental protocol, and Section~\ref{sec:conclusion} concludes the paper. Fig.~\ref{fig:overview} summarizes the proposed framework.

\begin{figure*}[!t]
\centering
\figph{\textwidth}{5.6cm}{\textbf{[Placeholder] Overall framework.}\\[2pt]
Left to right: measurement window $\mX_k$ and past-input window $\mU_k^{\mathrm{in}}$ $\rightarrow$ Transformer lifting $g_\theta$ and context map $C_\phi$ $\rightarrow$ lifted state $\vz_k$ and history context $\vc_k$ $\rightarrow$ scheduled matrices $\mA(\vc_k)$, $\mB(\vchi_k)$ with nominal $\mA_0$, $\mB_0$ and low-rank corrections $\rightarrow$ $\vz_{k+1}=\mA(\vc_k)\vz_k+\mB(\vchi_k)\vu_k$ and readout $\vx=C_x\vz$.\\[2pt]
Lower branch: nominal rollout $\rightarrow$ frozen $\bar\mA_i$, $\bar\mB_i$ $\rightarrow$ convex QP in measurement space $\rightarrow$ input sequence $U$.}
\caption{Overall framework of the history-conditioned scheduled predictor and the frozen-LTV controller. (Placeholder; figure to be drawn.)}
\label{fig:overview}
\end{figure*}

% =========================================================
\section{Preliminaries and Problem Formulation}
\label{sec:prelim}
% =========================================================

\subsection{Coordinates, controlled Koopman representation, and notation}

Consider a discrete-time controlled system
\begin{equation}
    \vs_{k+1}=f(\vs_k,\vu_k),\qquad \vx_k=h(\vs_k),
    \label{eq:physical_system}
\end{equation}
where $\vs_k\in\mathcal S$ is a physical state, $\vu_k\in\R^m$ is the input, and $\vx_k\in\R^n$ is the available measurement. For a fixed input, the controlled Koopman operator acts on an observable $\varphi$ according to
\begin{equation}
    (\mathcal K_{\vu}\varphi)(\vs)=\varphi(f(\vs,\vu)).
    \label{eq:koopman_operator}
\end{equation}
It is linear in the observable function, although the underlying state dynamics need not be linear. The finite-dimensional matrices used in this paper are a controlled prediction parameterization in selected lifted coordinates. They are not identified as an exact finite-dimensional truncation of the infinite-dimensional operator on the physical state.

For a vector $a$, let $\norm{a}_2$ denote the Euclidean norm; for a matrix, $\norm{M}_2$ denotes its spectral norm. For a positive semidefinite matrix $M$, write $\norm{a}_M^2=a^\top M a$. The symbol $P$ denotes the history length, whereas $P_x$ denotes a terminal cost matrix; these quantities are distinct.

\subsection{Finite histories and prediction/control tasks}

Let $P\ge2$. At time $k$, define the measurement and past-input windows
\begin{align}
    \mX_k&=\begin{bmatrix}\vx_{k-P+1}\\ \vdots\\ \vx_k\end{bmatrix}\in\R^{Pn},
    \label{eq:measurement_window}\\
    \mU_k^{\mathrm{in}}&=\begin{bmatrix}\vu_{k-P+1}\\ \vdots\\ \vu_{k-1}\end{bmatrix}\in\R^{(P-1)m}.
    \label{eq:history_windows}
\end{align}
The current input $\vu_k$ is not part of the context window; it enters the one-step update separately. Given $(\mX_k,\mU_k^{\mathrm{in}})$ and a future input sequence $\vu_{k:k+H-1}$, the prediction task is to estimate $\vx_{k+1:k+H}$. The control task is to choose a finite sequence $U=(\vu_0,\ldots,\vu_{N-1})$ that tracks a measurement reference while respecting the input set.

The lifted state is formed from the current measurement and a history encoding:
\begin{align}
    \vz_k&=g_\theta(\mX_k)=\begin{bmatrix}\vx_k\\ \psi_\theta(\mX_k)\end{bmatrix}\in\R^d,
    \label{eq:state_preserving_lifting}\\
    \vx_k&=C_x\vz_k,\qquad d=n+n_\psi,
    \label{eq:state_readout}
\end{align}
where $C_x=[I_n\;0]$. We distinguish $\vz_k^{\mathrm{enc}}=g_\theta(\mX_k)$, which uses measured data, from a recursively propagated $\vz_k^{\mathrm{prop}}$, whose measured coordinates are model outputs.

The model and the control analysis use a common-history convention. All compared rollouts start from the same measured window and past-input window and from the same $\vz_0=g_\theta(\mX_k)$. For a predicted index $j<0$, a history slot reads the corresponding measured past sample rather than a fictitious negative-time latent state. After each step, the oldest measurement slot is replaced by $C_x\vz_{i+1}$ and the oldest input slot by $\vu_i$. Fig.~\ref{fig:rollout} illustrates this update.

\begin{figure}[!t]
\centering
\figph{\columnwidth}{4.2cm}{\textbf{[Placeholder] Recursive window update.}\\[2pt]
Sliding windows $\mX_k$, $\mU_k^{\mathrm{in}}$: measured past samples (known) and predicted samples $C_x\vz_{i+1}$ (model output) entering the window step by step; oldest slot dropped; $\vz_i^{\mathrm{enc}}$ vs.\ $\vz_i^{\mathrm{prop}}$.}
\caption{Common-history convention and recursive window update. (Placeholder; figure to be drawn.)}
\label{fig:rollout}
\end{figure}

% =========================================================
\section{Scheduled Koopman Modeling}
\label{sec:model}
% =========================================================

\subsection{Model construction}

The intended architecture is a patch-based temporal Transformer encoder \cite{vaswani2017attention,nie2023time}. We impose the implementation constraint $p\mid P$: the history is partitioned into $P/p$ contiguous patches, each patch is linearly embedded, augmented with positional encoding, processed by a Transformer encoder, and pooled before projection to the hidden lifted coordinates. The current measurement is concatenated to this projection as in \eqref{eq:state_preserving_lifting}. Thus, the theoretical map $g_\theta$ abstracts the architecture while the measurement readout and the recursive propagation interface remain explicit. The context map $C_\phi$ is a separate feed-forward map of the flattened measurement and past-input windows. The bounds below depend on this architecture only through the stated Lipschitz and work-domain conditions.

The history context is
\begin{align}
    \vc_k&=C_\phi(\mX_k,\mU_k^{\mathrm{in}})\in\R^{d_c},
    \label{eq:history_context}\\
    \vchi_k&=C_x\vz_k\in\R^n.
    \label{eq:scheduling_variables}
\end{align}
The scheduled predictor is
\begin{equation}
    \vz_{k+1}=\mA(\vc_k)\vz_k+\mB(\vchi_k)\vu_k.
    \label{eq:scheduled_predictor}
\end{equation}
The map $\mA$ is scheduled by history, while $\mB$ is scheduled by the current measured component of the lifted state. This separation is a parameterization choice.

We use bounded low-rank corrections around nominal matrices,
\begin{align}
    \mA(\vc)&=\mA_0+\mU_A\operatorname{diag}(\gamma_A(\vc))\mV_A^\top,
    \label{eq:transition_schedule}\\
    \mB(\vchi)&=\mB_0+\mU_B\operatorname{diag}(\gamma_B(\vchi))\mV_B^\top.
    \label{eq:low_rank_schedules}
\end{align}
Here $r_A,r_B\in\mathbb N_{\ge1}$, $g_A,g_B\ge0$, $\mA_0\in\R^{d\times d}$, $\mB_0\in\R^{d\times m}$, $\mU_A,\mV_A\in\R^{d\times r_A}$, $\mU_B\in\R^{d\times r_B}$, and $\mV_B\in\R^{m\times r_B}$. The transition coefficients have the form
\begin{align}
    \gamma_{A,j}(\vc)&=g_A\pi_{A,j}(\vc)\tanh([W_A\vc+b_A]_j),
    \label{eq:A_modulation_coefficients}\\
    \pi_A(\vc)&=\operatorname{softmax}(W_A^\pi\vc+b_A^\pi),
    \label{eq:A_modulation}
\end{align}
where $W_A,W_A^\pi\in\R^{r_A\times d_c}$ and $b_A,b_A^\pi\in\R^{r_A}$. The input-modulation coefficients are defined analogously by
\begin{align}
    \gamma_{B,j}(\vchi)&=g_B\pi_{B,j}(\vchi)\tanh([W_B\vchi+b_B]_j),\\
    \pi_B(\vchi)&=\operatorname{softmax}(W_B^\pi\vchi+b_B^\pi),
\end{align}
where $W_B,W_B^\pi\in\R^{r_B\times n}$ and $b_B,b_B^\pi\in\R^{r_B}$. Throughout this subsection, the column-normalization condition
\begin{equation}
 \norm{u_{A,j}}_2,\norm{v_{A,j}}_2,\norm{u_{B,j}}_2,\norm{v_{B,j}}_2\le1
 \label{eq:factor_norm_assumption}
\end{equation}
is a mathematical assumption on the factors. The present manuscript does not prescribe whether training enforces it by projection, reparameterization, or another mechanism; this implementation choice is marked as \textbf{TODO(author)}. It is not a stability condition: these assumptions do not imply $\norm{\mA(\vc)}_2<1$ or $\norm{\mB(\vchi)}_2<1$.

\begin{proposition}[Uniform modulation bounds]
\label{prop:modulation_bound}
For every $\vc\in\R^{d_c}$ and $\vchi\in\R^n$,
\begin{equation}
    \begin{aligned}
    \norm{\mA(\vc)-\mA_0}_2&\le g_A,\\
    \norm{\mB(\vchi)-\mB_0}_2&\le g_B.
    \end{aligned}
    \label{eq:modulation_bound}
\end{equation}
\end{proposition}
\begin{proof}
Writing the first correction as $\sum_{j=1}^{r_A}\gamma_{A,j}u_{A,j}v_{A,j}^{\top}$ and using $\norm{uv^{\top}}_2=\norm u_2\norm v_2$ gives $\norm{\mA(\vc)-\mA_0}_2\le\sum_{j=1}^{r_A}|\gamma_{A,j}(\vc)|$. Since $g_A\ge0$, softmax weights are nonnegative and sum to one, and $|\tanh(\cdot)|<1$, $|\gamma_{A,j}(\vc)|\le g_A\pi_{A,j}(\vc)$, which proves the first bound. The argument for $\mB$ is identical.
\end{proof}

The schedule maps are also Lipschitz under finite modulation-layer weights. Define
\begin{align}
    L_A&=g_A\left(\norm{W_A}_2+\frac{\sqrt{r_A}}{2}\norm{W_A^\pi}_2\right),
    \label{eq:transition_Lipschitz_constant}\\
    L_B&=g_B\left(\norm{W_B}_2+\frac{\sqrt{r_B}}{2}\norm{W_B^\pi}_2\right).
    \label{eq:Lipschitz_constants}
\end{align}
Then
\begin{equation}
    \begin{aligned}
    \norm{\mA(\vc)-\mA(\vc')}_2&\le L_A\norm{\vc-\vc'}_2,\\
    \norm{\mB(\vchi)-\mB(\vchi')}_2&\le L_B\norm{\vchi-\vchi'}_2.
    \end{aligned}
    \label{eq:schedule_lipschitz}
\end{equation}
The factor $1/2$ is a global Euclidean Lipschitz constant of softmax; for $r=1$ softmax is constant and the displayed bound remains valid conservatively. The constants are for a fixed trained model with finite weights and may be conservative; they are not claimed to be tight.

\subsection{Training and recursive prediction}

To make the training definition explicit, let $\mathcal I_1$ be the set of indices for which the measured window $\mX_k$, past-input window $\mU_k^{\mathrm{in}}$, current input $\vu_k$, and next measurement $\vx_{k+1}$ are available. For $k\in\mathcal I_1$, define
\begin{align}
 \vz_k^{\mathrm{enc}}&=g_\theta(\mX_k),\\
 \hat\vz_{k+1|k}&=\mA\!\left(C_\phi(\mX_k,\mU_k^{\mathrm{in}})\right)\vz_k^{\mathrm{enc}}
 +\mB(C_x\vz_k^{\mathrm{enc}})\vu_k,\\
 \vx_{k+1|k}&=C_x\hat\vz_{k+1|k},\qquad
 \vz_{k+1}^{\mathrm{enc}}=g_\theta(\mX_{k+1}).
 \label{eq:training_one_step_objects}
\end{align}
The one-step objective used in this theoretical specification is
\begin{align}
 \mathcal L_x^{(1)}&=\frac{1}{|\mathcal I_1|}\sum_{k\in\mathcal I_1}
 \norm{\vx_{k+1|k}-\vx_{k+1}}_2^2,\\
 \mathcal L_z^{(1)}&=\frac{1}{|\mathcal I_1|}\sum_{k\in\mathcal I_1}
 \norm{\hat\vz_{k+1|k}-\vz_{k+1}^{\mathrm{enc}}}_2^2,\\
 \mathcal L^{(1)}&=\mathcal L_x^{(1)}+w_z\mathcal L_z^{(1)},
 \qquad w_z\ge0.
 \label{eq:joint_training_loss}
\end{align}
All quantities in the right-hand side are evaluated using measured windows; in particular, $\mU_k^{\mathrm{in}}$ is supplied to the separate context network $C_\phi$, not to the Transformer lifting map $g_\theta$. The sums may be replaced by explicitly specified normalized weighted sums when samples or horizons are weighted; no such weighting is assumed here. The present manuscript does not establish a stop-gradient policy for the re-encoded target $\vz_{k+1}^{\mathrm{enc}}$; this is \textbf{TODO(author)} and does not affect the theorems below.

If a multi-step loss is used, let $\mathcal I_H$ contain samples with $H$ future inputs and measurements. The recursively generated $\hat\vz_{k+i|k}$ uses the propagated measurement and input windows defined below, while the measured future windows may be used only to form targets. A concrete multi-step objective would be
\begin{equation}
 \mathcal L^{(H)}=\frac{1}{H|\mathcal I_H|}\sum_{k\in\mathcal I_H}\sum_{i=1}^H
 \left(\norm{C_x\hat\vz_{k+i|k}-\vx_{k+i}}_2^2
 +w_z\norm{\hat\vz_{k+i|k}-\vz_{k+i}^{\mathrm{enc}}}_2^2\right),
 \label{eq:multistep_training_loss}
\end{equation}
but whether this optional objective is used, and whether target re-encodings participate in backpropagation, remains \textbf{TODO(author)}. In all cases, a finite training MSE is not equivalent to the uniform residual certificate in Assumption~\ref{ass:a5}.

At deployment, the model is recursively propagated. Starting from $\vz_0=g_\theta(\mX_k)$, construct the predicted history from the model's own outputs and use
\begin{equation}
    \vz_{i+1}=\mA(\vc_i^{\mathrm{prop}})\vz_i+\mB(C_x\vz_i)\vu_i,
    \label{eq:recursive_prediction}
\end{equation}
where $\vc_i^{\mathrm{prop}}$ is computed from the recursively shifted measurement and input windows. No future measured state is used in this recursion.

\subsection{Finite-horizon prediction analysis}

Let $\vz_{k+i}^{\mathrm{enc}}$ be the encoder sequence formed from the true measurement windows and let $\hat\vz_i$ be the propagated sequence under the same input sequence. Define
\begin{equation}
    \ve_i=\vz_{k+i}^{\mathrm{enc}}-\hat\vz_i,
    \qquad \ve_0=0.
    \label{eq:prediction_error}
\end{equation}
Let $\mathcal D_{\rm win}\subset\R^{Pn+(P-1)m}$ be a domain for the concatenated pair $(\mX,\mU^{\rm in})$, let $\mathcal Z\subset\R^d$ and $\mathcal C\subset\R^{d_c}$ be domains for lifted states and context outputs, and let $\mathcal U\subset\R^m$ be the single-step input domain. The maps have interfaces $g_\theta:\R^{Pn}\to\R^d$ and $C_\phi:\mathcal D_{\rm win}\to\mathcal C$. For the fixed prediction comparison, assume that the encoded and propagated trajectories remain in common finite domains, that $C_\phi$ is Lipschitz with constant $L_\phi$, and that the one-step latent residual
\begin{equation}
    \vr_{k}=\vz_{k+1}^{\mathrm{enc}}-\mA(\vc_k)\vz_k^{\mathrm{enc}}-\mB(\vchi_k)\vu_k
    \label{eq:residual_definition}
\end{equation}
obeys $\norm{\vr_k}_2\le\delta$ on the considered domain and horizon. Let
\begin{align}
    M_A&=\norm{\mA_0}_2+g_A,
    \label{eq:transition_norm_bound}\\
    L&=M_A+L_BM_U+L_AL_\phi\sqrt{P}M_Z,
    \label{eq:prediction_L}
\end{align}
where $M_Z$ and $M_U$ are common state and input bounds.

The assumptions below are local to a fixed finite horizon.
\begin{assumption}[A1: common finite work domain]
\label{ass:a1}
For fixed initial measured histories, parameters, horizon $H$, and future input sequence, there are compact sets $\mathcal Z$ and $\mathcal C$ and a compact input set $\mathcal U$ such that the truth-encoded and propagated trajectories used in this prediction comparison remain in $\mathcal Z$, their context outputs remain in $\mathcal C$, and their joint windows remain in $\mathcal D_{\rm win}$. Let $M_Z=\sup_{z\in\mathcal Z}\norm{z}_2$ and $M_U=\sup_{u\in\mathcal U}\norm{u}_2$.
\end{assumption}
\begin{remark}[Optional lifting regularity]
If a separate bound is desired for a lifting residual, one may assume that $g_\theta$ is Lipschitz on a named window domain. No such constant is used in Proposition~\ref{prop:prediction_recurrence} or Theorem~\ref{thm:prediction_bound}, so this condition is not part of their hypotheses.
\end{remark}
\begin{assumption}[A3: common-window context regularity]
\label{ass:a3}
On $\mathcal D_{\rm win}$, $C_\phi$ is Lipschitz with one constant $L_\phi$ for every pair of joint windows:
\begin{equation}
\begin{aligned}
&\norm{C_\phi(W,V)-C_\phi(W',V')}_2\\
&\quad\le L_\phi\norm{(W,V)-(W',V')}_2.
\end{aligned}
\end{equation}
\end{assumption}
\begin{remark}[Optional uniform model-family condition]
A uniform bound on the four modulation-layer weights would be needed only for a theorem over a family of trained models. The present results concern one fixed trained model with finite weights, so no model-family assumption is used.
\end{remark}
\begin{assumption}[A5: uniform one-step residual]
\label{ass:a5}
For all encoded transitions $k+i$ with $i=0,\ldots,H-1$ in the common work domain, the residual in \eqref{eq:residual_definition} satisfies $\norm{\vr_{k+i}}_2\le\delta$ for one common $\delta\ge0$.
\end{assumption}
Theorem~\ref{thm:prediction_bound} uses A1, A3, and A5 together with the structural bounds and the finite weights of the fixed trained model. The freezing bound uses the corresponding shared-domain and Lipschitz conditions for all candidate inputs. The conditional cost theorem additionally requires one common feasible set and a uniform cost approximation over that set. These constants may be analytically certified on a specified domain, estimated from finite data, or currently unknown; a finite-data estimate is not by itself a uniform theorem certificate.

\begin{proposition}[One-step prediction recurrence]
\label{prop:prediction_recurrence}
For every $0\le i\le H-1$,
\begin{equation}
\begin{aligned}
\norm{\ve_{i+1}}_2\le{}& \delta+(M_A+L_BM_U)\norm{\ve_i}_2\\
 &+L_AL_\phi\sqrt{P}M_Z
  \max_{i-P+1\le j\le i}\norm{\ve_j}_2,
\end{aligned}
\label{eq:prediction_recurrence}
\end{equation}
where $\norm{C_x}_2=1$ under $C_x=[I_n\;0]$. For slots preceding the prediction start, set $\ve_j=0$ only as an index convention for the shared measured history; this does not define a negative-time latent state.
\end{proposition}

\begin{proof}
Subtract the propagated update from the encoder update and add and subtract the encoder-scheduled terms:
\begin{equation}
\begin{aligned}
\ve_{i+1}={}&\vr_{k+i}\\
&+\left[\mA(\vc_{k+i}^{\mathrm{enc}})\vz_{k+i}^{\mathrm{enc}}-\mA(\vc_i^{\mathrm{prop}})\hat\vz_i\right]\\
&+\left[\mB(\vchi_{k+i}^{\mathrm{enc}})-\mB(\vchi_i^{\mathrm{prop}})\right]\vu_{k+i}.
\end{aligned}
\end{equation}
The first term is bounded by A5. For the transition branch, use
\begin{align*}
 &\mA(\vc_{k+i}^{\mathrm{enc}})\vz_{k+i}^{\mathrm{enc}}-
 \mA(\vc_i^{\mathrm{prop}})\hat\vz_i\\
 &=\mA(\vc_{k+i}^{\mathrm{enc}})\ve_i+
 [\mA(\vc_{k+i}^{\mathrm{enc}})-\mA(\vc_i^{\mathrm{prop}})]\hat\vz_i.
\end{align*}
The first term is bounded by $M_A\norm{\ve_i}_2$, and the second by $L_A M_Z$ times the context difference. The input portions of the two windows are identical, while each predicted measurement slot differs by $C_x\ve_j$ and each pre-horizon slot is shared. Thus the joint-window difference is at most $\norm{C_x}_2\sqrt P\max_{i-P+1\le j\le i}\norm{\ve_j}_2$; under $C_x=[I_n\;0]$, $\norm{C_x}_2=1$, and A3 gives the displayed context term. For the input branch, \eqref{eq:schedule_lipschitz} gives $L_B\norm{C_x\ve_i}_2\norm{\vu_{k+i}}_2\le L_BM_U\norm{\ve_i}_2$. Combining the three bounds proves \eqref{eq:prediction_recurrence}.
\end{proof}

\begin{theorem}[Finite-horizon open-loop prediction bound]
\label{thm:prediction_bound}
Under the conditions above, for every $0\le i\le H$,
\begin{align}
    \norm{\ve_i}_2&\le W_i\le \delta\sum_{t=0}^{i-1}L^{i-1-t},
    \label{eq:prediction_bound}
    \\
    W_i&=\max_{0\le j\le i}\norm{\ve_j}_2.
    \label{eq:prediction_error_maximum}
\end{align}
Thus, for $L\ne1$, $\norm{\ve_H}_2\le\delta(L^H-1)/(L-1)$, and for $L=1$, $\norm{\ve_H}_2\le H\delta$.
\end{theorem}
\begin{proof}
The residual, state-propagation, schedule-difference, and measurement-dependent input terms give \eqref{eq:prediction_recurrence}. Since the recent-window maximum is bounded by $W_i$, one obtains $\norm{\ve_{i+1}}_2\le\delta+LW_i$. Let $S_0=0$ and $S_{i+1}=\delta+LS_i$. The sequence $S_i=\delta\sum_{t=0}^{i-1}L^{i-1-t}$ is nondecreasing for $L\ge0$, since $S_{i+1}-S_i=\delta L^i$. Induction in $W_{i+1}=\max(W_i,\norm{\ve_{i+1}}_2)$ yields $W_i\le S_i$.
\end{proof}

The bound is open-loop and finite-horizon, and it does not require $L<1$. The residual bound $\delta$ and the work domain must be certified for a concrete system. The corresponding measurement-space error satisfies
\begin{equation}
 \norm{\vx_{k+i}-C_x\hat\vz_i}_2
 =\norm{C_x\ve_i}_2\le\norm{C_x}_2\norm{\ve_i}_2.
 \label{eq:measurement_prediction_bound}
\end{equation}
For $\delta=0$ the bound is zero; for $H=1$, $P=2$, $L=0$, and $L=1$ it reduces respectively to the direct one-step, zero-gain, and linear-growth cases. No claim of long-horizon boundedness follows unless additional conditions are imposed.

% =========================================================
\section{Frozen-LTV Model Predictive Control}
\label{sec:control}
% =========================================================

\subsection{Complete scheduled problem and freezing construction}

At a control instant, let $\vz_0=\vz_k^{\mathrm{enc}}$ and let $U=(\vu_0,\ldots,\vu_{N-1})$. Assume the input set $\mathcal U\subset\R^m$ is nonempty, compact, and convex; in the standard QP case it is a box or another polytope. For the input-only formulation, define the common feasible set $\mathcal U_{\mathrm{feas}}=\mathcal U^N$. The complete scheduled model uses
\begin{equation}
    \vz_{i+1}=\mA(\vc_i)\vz_i+\mB(C_x\vz_i)\vu_i,
    \label{eq:complete_control_model}
\end{equation}
with history slots shifted according to Section~\ref{sec:prelim}. Given measurement references $\vx_i^{\mathrm{ref}}$, input references $\vu_i^{\mathrm{ref}}$, previous input $\vu_{-1}$, and positive semidefinite weights $Q_x,P_x,R,S$, define
\begin{equation}
\begin{aligned}
J_{\mathrm{full}}(U)= {}&\sum_{i=1}^{N-1}\norm{C_x\vz_i-\vx_i^{\mathrm{ref}}}_{Q_x}^2\\
 &+\norm{C_x\vz_N-\vx_N^{\mathrm{ref}}}_{P_x}^2\\
 &+\sum_{i=0}^{N-1}\norm{\vu_i-\vu_i^{\mathrm{ref}}}_{R}^2\\
 &+\sum_{i=0}^{N-1}\norm{\vu_i-\vu_{i-1}}_{S}^2.
\end{aligned}
\label{eq:full_cost}
\end{equation}
The tracking cost is expressed in measurement space; no reference is assigned to the hidden lifted coordinates. The complete problem is generally nonlinear and nonconvex because the scheduling variables depend on the decision-dependent trajectory.

Choose a feasible nominal input sequence $\bar U=(\bar\vu_0,\ldots,\bar\vu_{N-1})$ and initialize the nominal rollout with the same state and pre-horizon history, $\bar\vz_0=\vz_0$, $\bar X_0=\mX_k$, and $\bar U^{\mathrm{in}}_0=\mU_k^{\mathrm{in}}$. At each step set $\bar\vchi_i=C_x\bar\vz_i$ and $\bar\vc_i=C_\phi(\bar X_i,\bar U_i^{\mathrm{in}})$, then roll out the complete model once:
\begin{equation}
    \bar\vz_{i+1}=\mA(\bar\vc_i)\bar\vz_i+\mB(\bar\vchi_i)\bar\vu_i.
    \label{eq:nominal_rollout}
\end{equation}
Set
\begin{align}
    \bar\mA_i&=\mA(\bar\vc_i),
    \label{eq:frozen_transition_matrix}\\
    \bar\mB_i&=\mB(\bar\vchi_i),
    \label{eq:frozen_matrices}
\end{align}
and keep these matrices fixed while optimizing the candidate sequence $U$. The frozen trajectory is therefore
\begin{equation}
    \vz_{i+1}^{\mathrm{fr}}=\bar\mA_i\vz_i^{\mathrm{fr}}+\bar\mB_i\vu_i.
    \label{eq:frozen_model}
\end{equation}
The nominal rollout, complete trajectory, and frozen trajectory are distinct objects: they use different input sequences and/or different scheduling rules. Fig.~\ref{fig:freezing} illustrates the construction.

\begin{figure}[!t]
\centering
\figph{\columnwidth}{4.6cm}{\textbf{[Placeholder] Freezing construction.}\\[2pt]
Three trajectories over the horizon: nominal $\bar\vz_i$ (input $\bar U$), complete $\vz_i^{\mathrm{full}}$ (scheduling recomputed along the trajectory), frozen $\vz_i^{\mathrm{fr}}$ (matrices $\bar\mA_i$, $\bar\mB_i$ fixed). Gaps $\vg_i$ (nominal drift) and $\ve_i$ (complete vs.\ frozen).}
\caption{Nominal, complete, and frozen trajectories in the frozen-LTV construction. (Placeholder; figure to be drawn.)}
\label{fig:freezing}
\end{figure}

For the uniform freezing and cost statements below, fix $\bar U\in\mathcal U_{\mathrm{feas}}$ and assume the following control comparison condition: for every candidate $U\in\mathcal U_{\mathrm{feas}}$, the full, frozen, and nominal trajectories remain in shared compact lifted-state and context-output domains, all joint windows remain in $\mathcal D_{\rm win}$, and the same constants $M_Z$, $M_U$, and $L_\phi$ apply. The input-window domain is $\mathcal D_{\rm win}$ and the context-output domain is $\mathcal C$; they are distinct. This condition is a hypothesis on the whole feasible input set, not an additional state constraint in the QP. The online policy used to select $\bar U$ is \textbf{TODO(author)}; the theory below is conditional on the chosen feasible nominal sequence and does not prescribe warm starts, shifted previous solutions, reference inputs, or iterative re-freezing.

\paragraph{One receding-horizon cycle.} At each real sampling time, update the measured windows, compute $\vz_0=g_\theta(\mX_k)$, obtain a feasible nominal $\bar U$ according to the author-confirmed policy, roll out the complete model once, freeze $\bar\mA_i,\bar\mB_i$, assemble and solve one convex QP, and apply its first input. At the next sampling time, rebuild the histories and repeat. If iterative re-freezing is used in an implementation, it is an additional algorithmic procedure; the fixed-nominal results here are not a convergence or closed-loop guarantee.

\subsection{Prediction matrices and convex quadratic program}

Stack the frozen states and inputs as
\begin{align}
    \mZ&=[(\vz_1^{\mathrm{fr}})^\top,\ldots,(\vz_N^{\mathrm{fr}})^\top]^\top,
    \label{eq:stacked_frozen_states}\\
    \mU&=[\vu_0^\top,\ldots,\vu_{N-1}^\top]^\top.
    \label{eq:stacked_inputs}
\end{align}
The frozen model gives
\begin{equation}
    \mZ=\mF+\mG\mU,
    \label{eq:frozen_prediction_matrices}
\end{equation}
where the block rows of $\mF$ and $\mG$ are denoted by $F_i$ and $G_{i,t}$, respectively:
\begin{align}
    F_i&=\bar\mA_{i-1}\cdots\bar\mA_0\vz_0,
    \label{eq:F_blocks}\\
    G_{i,t}&=\begin{cases}
    \bar\mA_{i-1}\cdots\bar\mA_{t+1}\bar\mB_t,&t<i,\\
    0,&t\ge i.
    \end{cases}
    \label{eq:F_G_blocks}
\end{align}
The products are time ordered; no commutativity of the matrices is assumed. Let $\bar C_x=I_N\otimes C_x$, $Y^{\mathrm{ref}}=[(\vx_1^{\mathrm{ref}})^\top,\ldots,(\vx_N^{\mathrm{ref}})^\top]^\top$, and $\bar Q_x=\operatorname{blkdiag}(Q_x,\ldots,Q_x,P_x)$. Define $\bar R$ and $\bar S$ as the corresponding block-diagonal input matrices, and let $\hat D$ satisfy $(\hat D\mU)_0=\vu_0$ and $(\hat D\mU)_i=\vu_i-\vu_{i-1}$ for $i\ge1$. With the input-offset vector $\hat{\mathbf c}=[-\vu_{-1}^\top,0,\ldots,0]^\top$, the frozen cost is
\begin{equation}
\begin{aligned}
J_{\mathrm{fr}}(U;\bar U)={}&\norm{\bar C_x\mZ-Y^{\mathrm{ref}}}_{\bar Q_x}^2\\
&+\norm{\mU-\mU^{\mathrm{ref}}}_{\bar R}^2\\
&+\norm{\hat D\mU+\hat{\mathbf c}}_{\bar S}^2.
\end{aligned}
\label{eq:frozen_cost}
\end{equation}
Here $\mF\in\R^{Nd}$, $\mG\in\R^{Nd\times Nm}$, $\bar C_x\in\R^{Nn\times Nd}$, and $\hat D\in\R^{Nm\times Nm}$. The first row of $\hat D$ is $[I_m\;0\;\cdots\;0]$ and row $i\ge1$ has $-I_m$ in block column $i-1$ and $I_m$ in block column $i$; thus $\hat{\mathbf c}=[-\vu_{-1}^{\top},0,\ldots,0]^{\top}\in\R^{Nm}$. Empty matrix products in \eqref{eq:F_G_blocks} are identity matrices. For a polyhedral input set $\mathcal U=\{u:A_u u\le b_u\}$, the stacked constraint is $(I_N\otimes A_u)\mU\le\mathbf 1_N\otimes b_u$.

Substitution of \eqref{eq:frozen_prediction_matrices} gives
\begin{equation}
    J_{\mathrm{fr}}(U;\bar U)=\frac12\mU^\top H\mU+q^\top\mU+\mathrm{const},
    \label{eq:quadratic_template}
\end{equation}
where
\begin{align}
H&=2\left(\begin{aligned}
 &\mG^\top\bar C_x^\top\bar Q_x\bar C_x\mG\\
 &+\bar R+\hat D^\top\bar S\hat D
\end{aligned}\right),
\label{eq:QP_Hessian}\\
 q&=2f,\qquad f=\begin{aligned}
 &\mG^\top\bar C_x^\top\bar Q_x(\bar C_x\mF-Y^{\mathrm{ref}})\\
 &-\bar R\mU^{\mathrm{ref}}+\hat D^\top\bar S\hat{\mathbf c}.
\end{aligned}
\label{eq:QP_coefficients}
\end{align}
Every term in $H$ is positive semidefinite, so $H\succeq0$ and the objective is convex. With a box or other polyhedral input set, the frozen problem is a standard convex QP. Strict convexity and uniqueness require an additional positive-definiteness condition, for example $\bar R\succ0$. When $N=1$, the output term contains only the terminal weight $P_x$ and no intermediate $Q_x$ term.

\subsection{Freezing error and conditional suboptimality}

Fix a candidate $U$ and nominal $\bar U$. Let $\vz_i^{\mathrm{full}}$ be the complete scheduled trajectory, $\vz_i^{\mathrm{fr}}$ the frozen trajectory, and $\bar\vz_i$ the nominal trajectory. For $i=0,\ldots,N-1$, let $U_i^{\mathrm{in}}$ and $\bar U_i^{\mathrm{in}}$ denote the length-$(P-1)$ past-input windows assembled from the common pre-horizon inputs and from $\vu_0,\ldots,\vu_{i-1}$ or $\bar\vu_0,\ldots,\bar\vu_{i-1}$, respectively. Define $\ve_i=\vz_i^{\mathrm{full}}-\vz_i^{\mathrm{fr}}$ and $\vg_i=\vz_i^{\mathrm{fr}}-\bar\vz_i$. Let
\begin{align}
    \delta_u&=\max_{0\le i<N}\norm{\vu_i-\bar\vu_i}_2,
    \label{eq:input_deviation}\\
    \Delta_u&=\max_{0\le i<N}\norm{U_i^{\mathrm{in}}-\bar U_i^{\mathrm{in}}}_2.
    \label{eq:input_window_deviation}
\end{align}
Then $\Delta_u\le\sqrt{P-1}\delta_u$ because the past input windows share their initial entries and differ only in shifted future inputs. With $M_B=\norm{\mB_0}_2+g_B$, the nominal drift satisfies
\begin{equation}
    \norm{\vg_i}_2\le G_i:=M_B\delta_u\sum_{t=0}^{i-1}M_A^{i-1-t}.
    \label{eq:nominal_drift}
\end{equation}

The three trajectories have a useful exact first-step identity. Because their initial state and history are common, $\bar\mA_0=\mA(\vc_0^{\mathrm{full}})$ and $\bar\mB_0=\mB(\vchi_0^{\mathrm{full}})$ for every candidate $U$. Consequently,
\begin{equation}
 \vz_1^{\mathrm{full}}=\vz_1^{\mathrm{fr}},\qquad \ve_0=\ve_1=0,
 \label{eq:first_step_freezing_exact}
\end{equation}
even when $U\ne\bar U$. This exactness is not a closed-loop statement.

For a less conservative candidate-specific bound, define $a_i=\norm{\vu_i-\bar\vu_i}_2$, $D_i=\norm{U_i^{\mathrm{in}}-\bar U_i^{\mathrm{in}}}_2$, and
\begin{equation}
 G_0^{\mathrm{st}}=0,\qquad G_{i+1}^{\mathrm{st}}=M_A G_i^{\mathrm{st}}+M_Ba_i,
 \label{eq:stage_nominal_drift}
\end{equation}
so that $\norm{\vg_i}_2\le G_i^{\mathrm{st}}$. Define the history envelope
\begin{equation}
 \widehat G_i:=\max_{\max\{0,i-P+1\}\le j\le i}G_j^{\mathrm{st}},
 \label{eq:stage_history_drift}
\end{equation}
where pre-horizon slots have drift zero. Let $E_0=0$ and recursively set
\begin{equation}
\begin{aligned}
 d_i&=L_AL_\phi M_ZD_i
 +(L_AL_\phi\sqrt P M_Z+L_BM_U)\widehat G_i,\\
 E_{i+1}&=\max\{E_i,LE_i+d_i\}.
\end{aligned}
\label{eq:stage_freezing_bound}
\end{equation}
Under the shared-domain conditions above, $\norm{\ve_i}_2\le E_i$. In particular, $D_0=G_0^{\mathrm{st}}=\widehat G_0=d_0=E_1=0$, which agrees with \eqref{eq:first_step_freezing_exact}. Neither the pointwise input-window deviations nor the stage drifts need be monotone, so no geometric-sum simplification of \eqref{eq:stage_freezing_bound} is made.

To make the bound uniform over the complete feasible set, define the fixed-nominal suprema
\begin{align}
\bar\delta_u&:=\sup_{U\in\mathcal U_{\mathrm{feas}}}\max_{0\le i<N}\norm{\vu_i-\bar\vu_i}_2,
\label{eq:uniform_pointwise_input_deviation}\\
\bar\Delta_u&:=\sup_{U\in\mathcal U_{\mathrm{feas}}}\max_{0\le i<N}\norm{U_i^{\mathrm{in}}-\bar U_i^{\mathrm{in}}}_2.
\label{eq:uniform_input_deviation}
\end{align}
Compactness of $\mathcal U_{\mathrm{feas}}$ makes these quantities finite, and $\bar\Delta_u\le\sqrt{P-1}\bar\delta_u$. Set
\begin{align}
\bar G_N&:=M_B\bar\delta_u\sum_{t=0}^{N-1}M_A^t,
\label{eq:uniform_nominal_drift}\\
\bar\delta_{\mathrm{frz}}&:=\begin{aligned}
 &L_AL_\phi M_Z\bar\Delta_u\\
 &+\left(L_AL_\phi\sqrt{P}M_Z+L_BM_U\right)\bar G_N,
\end{aligned}
\label{eq:uniform_freezing_driver}
\end{align}
and, for $0\le i\le N$,
\begin{equation}
\bar K_i:=\bar\delta_{\mathrm{frz}}\sum_{t=0}^{i-1}L^{i-1-t}.
\label{eq:uniform_freezing_bound}
\end{equation}
For every $U\in\mathcal U_{\mathrm{feas}}$, the candidate-specific quantities satisfy $\delta_u\le\bar\delta_u$, $\Delta_u\le\bar\Delta_u$, $G_N\le\bar G_N$, and hence $\norm{\ve_i}_2\le\bar K_i$.

For a common work domain, define
\begin{equation}
    \begin{aligned}
    \delta_{\mathrm{frz}}={}&L_AL_\phi M_Z\Delta_u\\
    &+\left(L_AL_\phi\sqrt{P}M_Z+L_BM_U\right)G_N.
    \end{aligned}
    \label{eq:freezing_driver}
\end{equation}
With $L$ from \eqref{eq:prediction_L}, the freezing error obeys
\begin{equation}
    \norm{\ve_i}_2\le K_i:=\delta_{\mathrm{frz}}
    \sum_{t=0}^{i-1}L^{i-1-t},\qquad 0\le i\le N.
    \label{eq:freezing_bound}
\end{equation}
This is a same-input, finite-horizon open-loop comparison.

The stage-dependent bound \eqref{eq:stage_freezing_bound} is the sharper candidate-specific certificate. The simpler $K_i$ bound above is retained as a uniform, conservative certificate because it uses horizon-wide envelopes.

To express the cost difference uniformly, let $M_Y=\norm{C_x}_2M_Z$, $M_{\mathrm{ref}}=\max_{1\le i\le N}\norm{\vx_i^{\mathrm{ref}}}_2$, and $\bar K_i^x=\norm{C_x}_2\bar K_i$. A uniform sufficient bound is
\begin{equation}
\begin{aligned}
\varepsilon_J={}&2(M_Y+M_{\mathrm{ref}})
 \left[\begin{aligned}
 &\norm{Q_x}_2\sum_{i=1}^{N-1}\bar K_i^x\\
 &+\norm{P_x}_2\bar K_N^x
 \end{aligned}\right]\\
&+\norm{Q_x}_2\sum_{i=1}^{N-1}(\bar K_i^x)^2\\
&+\norm{P_x}_2(\bar K_N^x)^2.
\end{aligned}
\label{eq:epsilon_cost}
\end{equation}
It follows from the identity for the difference of two squared weighted norms and from the fact that the input and input-increment terms in \eqref{eq:full_cost} and \eqref{eq:frozen_cost} are identical. More explicitly, let $y_i^{\mathrm{full}}=C_x\vz_i^{\mathrm{full}}$, $y_i^{\mathrm{fr}}=C_x\vz_i^{\mathrm{fr}}$, and $d_i=y_i^{\mathrm{full}}-y_i^{\mathrm{fr}}$. The uniform freezing bound gives $\norm{d_i}_2\le\bar K_i^x$. With $M_Y=\norm{C_x}_2M_Z$ and $M_{\mathrm{ref}}=\max_{1\le i\le N}\norm{\vx_i^{\mathrm{ref}}}_2$, A1 gives $\norm{y_i^{\mathrm{full}}}_2\le M_Y$ and hence $\norm{y_i^{\mathrm{fr}}}_2\le M_Y+\bar K_i^x$. Therefore, for $i<N$,
\begin{equation}
\begin{aligned}
&\left|\norm{y_i^{\mathrm{full}}-\vx_i^{\mathrm{ref}}}_{Q_x}^2-
\norm{y_i^{\mathrm{fr}}-\vx_i^{\mathrm{ref}}}_{Q_x}^2\right|\\
&\quad\le\norm{Q_x}_2\bar K_i^x(2M_Y+\bar K_i^x+2M_{\mathrm{ref}}),
\end{aligned}
\end{equation}
and the same identity with $P_x$ applies at $i=N$. Summing these bounds and using the common input terms yields \eqref{eq:epsilon_cost}. Because \eqref{eq:uniform_input_deviation} takes the supremum over the entire feasible set and A1/A3 use the same domain constants there, this calculation proves the following proposition.

\begin{proposition}[Uniform objective approximation certificate]
\label{prop:uniform_cost_certificate}
Under the shared-domain conditions above,
\begin{equation}
 \sup_{U\in\mathcal U_{\mathrm{feas}}}
 \left|J_{\mathrm{full}}(U)-J_{\mathrm{fr}}(U;\bar U)\right|\le\varepsilon_J.
 \label{eq:uniform_cost_certificate}
\end{equation}
\end{proposition}
\begin{remark}[Abstract uniform cost assumption]
\label{ass:cost_approximation}
The proof above supplies a sufficient certificate for the following abstract condition. For the fixed nominal sequence, the complete and frozen costs obey
\begin{equation}
    \sup_{U\in\mathcal U_{\mathrm{feas}}}
    |J_{\mathrm{full}}(U)-J_{\mathrm{fr}}(U;\bar U)|\le\varepsilon_J.
\end{equation}
\end{remark}

\begin{theorem}[Conditional cost suboptimality]
\label{thm:suboptimality}
Suppose the complete and frozen problems have the same nonempty feasible input set, and let $U^{\mathrm{fr}}$ and $U^\star$ be global minimizers of the frozen and complete costs, respectively. Under the abstract uniform cost condition in Remark~\ref{ass:cost_approximation},
\begin{equation}
    0\le J_{\mathrm{full}}(U^{\mathrm{fr}})-J_{\mathrm{full}}(U^\star)
    \le2\varepsilon_J.
    \label{eq:suboptimality_bound}
\end{equation}
\end{theorem}
\begin{proof}
The lower bound follows from global optimality of $U^\star$. For the upper bound,
\begin{align*}
J_{\mathrm{full}}(U^{\mathrm{fr}})
&\le J_{\mathrm{fr}}(U^{\mathrm{fr}};\bar U)+\varepsilon_J\\
&\le J_{\mathrm{fr}}(U^\star;\bar U)+\varepsilon_J\\
&\le J_{\mathrm{full}}(U^\star)+2\varepsilon_J.
\end{align*}
\end{proof}

The theorem applies only to global minimizers on the same fixed nonempty feasible input set. Existence may follow from continuity of the costs and compactness of $\mathcal U_{\mathrm{feas}}$, but the theorem does not imply that the nonconvex complete global minimizer can be computed efficiently. If a solver returns $\widetilde U$ satisfying $J_{\mathrm{fr}}(\widetilde U;\bar U)\le J_{\mathrm{fr}}(U^{\mathrm{fr}};\bar U)+\eta$, the same argument gives the separately conditional bound $J_{\mathrm{full}}(\widetilde U)-J_{\mathrm{full}}(U^\star)\le2\varepsilon_J+\eta$. If $\varepsilon_J$ is large over the full feasible set, the result remains mathematically valid but may have little performance-discrimination value. The comparison is per sampling instant with a fixed nominal sequence; it does not establish closed-loop stability, recursive feasibility, robust constraint satisfaction, or a multi-step closed-loop performance guarantee.

% =========================================================
\section{Experiments}
\label{sec:experiments}
% =========================================================

This section is intentionally a manuscript placeholder. Experimental results, figures, and numerical claims are omitted in the present draft and must be inserted only after the corresponding data, code version, and evidence have been independently checked. No statement in Sections~\ref{sec:introduction}--\ref{sec:control} relies on an unreported experimental result.

\subsection{Data and protocol}

The final version will specify the data source, sampling interval, state and input coordinates, history length, train/validation/test split, normalization, input constraints, and random seeds. The protocol must distinguish measurement coordinates from lifted coordinates and must report prediction horizon in both steps and physical time when the sampling interval is relevant.

\subsection{Modeling evaluation}

The final version will report one-step and multi-step measurement-space prediction errors, together with the joint training objective in \eqref{eq:joint_training_loss}. Any comparison among architectures or loss choices requires a fixed data split, fixed evaluation horizons, and an explicitly stated scope. The present theory does not claim that the lifted-state term improves empirical performance.

\subsection{Control evaluation}

The final version will describe the nominal-sequence policy, horizon, weights, input constraints, solver status, and whether each result is an offline simulation or a physical experiment. Offline closed-loop simulation must not be described as hardware validation. Timing measurements, if later reported, must be measured independently rather than inferred from the convex-QP formulation.

% =========================================================
\section{Conclusion}
\label{sec:conclusion}
% =========================================================

This paper formulated a history-conditioned scheduled predictor with a temporal lifting, a bounded low-rank history-conditioned transition, and a measurement-dependent input modulation. The scheduled form preserves linear propagation when the scheduling signals are fixed, while its recursive use remains generally nonlinear. The finite-horizon prediction result makes the history-window coupling explicit and gives a conditional open-loop bound under common-history, work-domain, Lipschitz, and residual conditions.

For control, freezing the scheduling matrices along a nominal rollout produces an LTV prediction model and a convex QP with measurement-space tracking costs. The freezing analysis separates nominal drift from the complete-versus-frozen trajectory error. Under a uniform cost approximation and a common feasible set, the frozen global optimizer satisfies a conditional $2\varepsilon_J$ cost bound relative to the global optimizer of the complete scheduled problem.

The results are finite-horizon and open-loop; the uniform work-domain and residual conditions require certification for a concrete system.

\appendices

\section{Proof Details}
\label{app:proofs}

\subsection*{Scheduling bounds}
For completeness, the softmax Jacobian is $J(\pi)=\operatorname{diag}(\pi)-\pi\pi^\top$. It is positive semidefinite and, by the variance bound, satisfies $\norm{J(\pi)}_2\le1/2$. Integrating the Jacobian along the line segment between two logits gives $\norm{\operatorname{softmax}(a)-\operatorname{softmax}(b)}_2\le\frac12\norm{a-b}_2$. Combining this inequality with the $\tanh$ Lipschitz property, $\norm{\cdot}_1\le\sqrt r\norm{\cdot}_2$, and the rank-one factor normalization gives \eqref{eq:schedule_lipschitz}.

\subsection*{Frozen error bound}
For a fixed candidate $U$ and nominal $\bar U$, the three trajectories are distinct: $\bar\vz_i$ uses $\bar U$ and recomputed schedules, $\vz_i^{\mathrm{fr}}$ uses $U$ and fixed nominal matrices, and $\vz_i^{\mathrm{full}}$ uses $U$ and recomputed schedules. Their differences satisfy
\begin{equation}
 \vg_{i+1}=\bar\mA_i\vg_i+\bar\mB_i(\vu_i-\bar\vu_i),
 \qquad
 \ve_{i+1}=\bar\mA_i\ve_i+[\mA(\vc_i^{\mathrm{full}})-\bar\mA_i]\vz_i^{\mathrm{full}}+[\mB(\vchi_i^{\mathrm{full}})-\bar\mB_i]\vu_i.
 \label{eq:appendix_three_trajectory_recursions}
\end{equation}
Thus $\norm{\vg_i}_2\le G_i$ in \eqref{eq:nominal_drift}. The full and nominal measurement histories differ slotwise by $C_x(\ve_j+\vg_j)$ after the prediction start and are identical before it, so
\begin{equation*}
 \norm{\mX_i^{\mathrm{full}}-\bar\mX_i}_2
 \le\norm{C_x}_2\sqrt P(W_i+G_N).
\end{equation*}
The input-window difference is at most $\Delta_u$. A3 therefore bounds the context difference by $L_\phi[\sqrt P(W_i+G_N)+\Delta_u]$. The current measured-state difference is at most $\norm{C_x}_2(W_i+G_N)$, so the schedule Lipschitz bound gives
\begin{equation*}
 \norm{\ve_{i+1}}_2\le\delta_{\mathrm{frz}}+LW_i,
\end{equation*}
with $\delta_{\mathrm{frz}}$ from \eqref{eq:freezing_driver}; the original manuscript uses $\norm{C_x}_2=1$. The nondecreasing comparison sequence then proves \eqref{eq:freezing_bound}.

For the stage-dependent refinement, $D_0=G_0^{\mathrm{st}}=0$, hence the two full schedules at stage zero equal the nominal schedules and $\ve_1=0$. At stage $i$, the same history argument with $D_i$ and $G_i^{\mathrm{st}}$ gives $\norm{\ve_{i+1}}_2\le LE_i+d_i$. Taking the maximum with $E_i$ proves by induction that $\norm{\ve_i}_2\le E_i$. No monotonicity of $D_i$ is used.

\subsection*{Quadratic program}
Let $\mG_x=\bar C_x\mG$ and $\mW=\bar C_x\mF-Y^{\mathrm{ref}}$. The output term expands as
\begin{equation}
    \begin{aligned}
    &(\mW+\mG_x\mU)^\top\bar Q_x(\mW+\mG_x\mU)\\
    &\quad=\mU^\top\mG_x^\top\bar Q_x\mG_x\mU\\
    &\qquad+2(\mG_x^\top\bar Q_x\mW)^\top\mU+\mathrm{const}.
    \end{aligned}
\end{equation}
The input and input-increment terms expand analogously. Their sum yields \eqref{eq:QP_Hessian} and $q=2f$ in \eqref{eq:QP_coefficients}. Positive semidefiniteness follows because each quadratic coefficient has the form $T^\top MT$ with $M\succeq0$.



\begin{thebibliography}{99}

\bibitem{wu2026multistep}
L. Wu, W. G. Y. Tan, L. Zhou, R. D. Braatz, and J. Drgo\v{n}a, ``Least-squares multi-step Koopman operator learning for model predictive control,'' arXiv:2601.11901, 2026.

\bibitem{ljung1999system}
L. Ljung, \emph{System Identification: Theory for the User}, 2nd ed. Upper Saddle River, NJ, USA: Prentice Hall PTR, 1999.

\bibitem{vanoverschee1994n4sid}
P. Van Overschee and B. De Moor, ``N4SID: Subspace algorithms for the identification of combined deterministic-stochastic systems,'' \emph{Automatica}, vol. 30, no. 1, pp. 75--93, 1994.

\bibitem{nelles2020nonlinear}
O. Nelles, \emph{Nonlinear System Identification: From Classical Approaches to Neural Networks, Fuzzy Models, and Gaussian Processes}, 2nd ed. Cham, Switzerland: Springer, 2020.

\bibitem{forgione2021dynonet}
M. Forgione and D. Piga, ``dynoNet: A neural network architecture for learning dynamical systems,'' \emph{Int. J. Adapt. Control Signal Process.}, vol. 35, no. 4, pp. 612--626, 2021.

\bibitem{koopman1931hamiltonian}
B. O. Koopman, ``Hamiltonian systems and transformation in Hilbert space,'' \emph{Proc. Natl. Acad. Sci. USA}, vol. 17, no. 5, pp. 315--318, 1931.

\bibitem{mezic2005spectral}
I. Mezi\'c, ``Spectral properties of dynamical systems, model reduction and decompositions,'' \emph{Nonlinear Dyn.}, vol. 41, no. 1--3, pp. 309--325, 2005.

\bibitem{williams2015data}
M. O. Williams, I. G. Kevrekidis, and C. W. Rowley, ``A data-driven approximation of the Koopman operator: Extending dynamic mode decomposition,'' \emph{J. Nonlinear Sci.}, vol. 25, no. 6, pp. 1307--1346, 2015.

\bibitem{proctor2016dmdc}
J. L. Proctor, S. L. Brunton, and J. N. Kutz, ``Dynamic mode decomposition with control,'' \emph{SIAM J. Appl. Dyn. Syst.}, vol. 15, no. 1, pp. 142--161, 2016.

\bibitem{proctor2018generalizing}
J. L. Proctor, S. L. Brunton, and J. N. Kutz, ``Generalizing Koopman theory to allow for inputs and control,'' \emph{SIAM J. Appl. Dyn. Syst.}, vol. 17, no. 1, pp. 909--930, 2018.

\bibitem{korda2018linear}
M. Korda and I. Mezi\'c, ``Linear predictors for nonlinear dynamical systems: Koopman operator meets model predictive control,'' \emph{Automatica}, vol. 93, pp. 149--160, 2018.

\bibitem{brunton2022modern}
S. L. Brunton, M. Budi\v{s}i\'c, E. Kaiser, and J. N. Kutz, ``Modern Koopman theory for dynamical systems,'' \emph{SIAM Rev.}, vol. 64, no. 2, pp. 229--340, 2022.

\bibitem{otto2021koopman}
S. E. Otto and C. W. Rowley, ``Koopman operators for estimation and control of dynamical systems,'' \emph{Annu. Rev. Control Robot. Auton. Syst.}, vol. 4, no. 1, pp. 59--87, 2021.

\bibitem{duran2025control}
J. F. Dur\'an-Sig\"uenza, L. I. Minchala, L. E. Garza-Casta\~n\'on, and H. Zhang, ``Control based on the Koopman operator: A comprehensive review,'' \emph{J. Franklin Inst.}, vol. 362, no. 18, art. 108256, 2025.

\bibitem{strasser2026overview}
R. Str\"asser, K. Worthmann, I. Mezi\'c, J. Berberich, M. Schaller, and F. Allg\"ower, ``An overview of Koopman-based control: From error bounds to closed-loop guarantees,'' \emph{Annu. Rev. Control}, vol. 61, art. 101035, 2026.

\bibitem{lusch2018deep}
B. Lusch, J. N. Kutz, and S. L. Brunton, ``Deep learning for universal linear embeddings of nonlinear dynamics,'' \emph{Nat. Commun.}, vol. 9, no. 1, p. 4950, 2018.

\bibitem{takeishi2017koopman}
N. Takeishi, Y. Kawahara, and T. Yairi, ``Learning Koopman invariant subspaces for dynamic mode decomposition,'' in \emph{Proc. NeurIPS}, 2017, pp. 1130--1140.

\bibitem{otto2019linearly}
S. E. Otto and C. W. Rowley, ``Linearly-recurrent autoencoder networks for learning dynamics,'' \emph{SIAM J. Appl. Dyn. Syst.}, vol. 18, no. 1, pp. 558--593, 2019.

\bibitem{han2020deep}
Y. Han, W. Hao, and U. Vaidya, ``Deep learning of Koopman representation for control,'' in \emph{Proc. IEEE Conf. Decision Control (CDC)}, 2020, pp. 1890--1895.

\bibitem{feizi2025deepkoopman}
N. Feizi, F. C. Pedrosa, J. Jayender, and R. V. Patel, ``Deep Koopman approach for nonlinear dynamics and control of tendon-driven continuum robots,'' \emph{IEEE Robot. Autom. Lett.}, vol. 10, no. 3, pp. 2926--2933, 2025.

\bibitem{brunton2017chaos}
S. L. Brunton, B. W. Brunton, J. L. Proctor, E. Kaiser, and J. N. Kutz, ``Chaos as an intermittently forced linear system,'' \emph{Nat. Commun.}, vol. 8, art. 19, 2017.

\bibitem{kamb2020timedelay}
M. Kamb, E. Kaiser, S. L. Brunton, and J. N. Kutz, ``Time-delay observables for Koopman: Theory and applications,'' \emph{SIAM J. Appl. Dyn. Syst.}, vol. 19, no. 2, pp. 886--917, 2020.

\bibitem{arbabi2017ergodic}
H. Arbabi and I. Mezi\'c, ``Ergodic theory, dynamic mode decomposition, and computation of spectral properties of the Koopman operator,'' \emph{SIAM J. Appl. Dyn. Syst.}, vol. 16, no. 4, pp. 2096--2126, 2017.

\bibitem{das2019delay}
S. Das and D. Giannakis, ``Delay-coordinate maps and the spectra of Koopman operators,'' \emph{J. Stat. Phys.}, vol. 175, no. 6, pp. 1107--1145, 2019.

\bibitem{shenoy2025temporal}
S. Shenoy, B. Sharan, and H. Werner, ``Temporal autoencoder for identification and predictive control of nonlinear dynamics based on Koopman operator theory,'' in \emph{Proc. American Control Conf. (ACC)}, 2025, pp. 1795--1800.

\bibitem{han2024robust}
M. Han, Z. Li, Xiang Yin, and Xunyuan Yin, ``Robust learning and control of time-delay nonlinear systems with deep recurrent Koopman operators,'' \emph{IEEE Trans. Ind. Informat.}, vol. 20, no. 3, pp. 4675--4684, Mar. 2024.

\bibitem{forootani2025deepkoopformer}
A. Forootani, M. Khosravi, and M. Barati, ``DeepKoopFormer: A Koopman enhanced Transformer based architecture for time series forecasting,'' arXiv:2508.02616, 2025.

\bibitem{forootani2026spectral}
A. Forootani and R. Iervolino, ``Learnable Koopman-enhanced Transformer-based time series forecasting with spectral control,'' arXiv:2602.02592, 2026.

\bibitem{li2025mamko}
Z. Li, M. Han, and X. Yin, ``MamKO: Mamba-based Koopman operator for modeling and predictive control,'' in \emph{Proc. ICLR}, 2025.

\bibitem{liu2023koopa}
Y. Liu, C. Li, J. Wang, and M. Long, ``Koopa: Learning non-stationary time series dynamics with Koopman predictors,'' in \emph{Proc. NeurIPS}, 2023.

\bibitem{rugh2000gain}
W. J. Rugh and J. S. Shamma, ``Research on gain scheduling,'' \emph{Automatica}, vol. 36, no. 10, pp. 1401--1425, 2000.

\bibitem{leith1998velocity}
D. J. Leith and W. E. Leithead, ``Gain-scheduled and nonlinear systems: Dynamic analysis by velocity-based linearization families,'' \emph{Int. J. Control}, vol. 70, no. 2, pp. 289--317, 1998.

\bibitem{cisneros2016qlmpc}
P. S. G. Cisneros, S. Voss, and H. Werner, ``Efficient nonlinear model predictive control via quasi-LPV representation,'' in \emph{Proc. 55th IEEE Conf. Decision Control (CDC)}, 2016, pp. 3216--3221.

\bibitem{cisneros2020qlpv}
P. S. G. Cisneros and H. Werner, ``Nonlinear model predictive control for models in quasi-linear parameter varying form,'' \emph{Int. J. Robust Nonlinear Control}, vol. 30, no. 10, pp. 3945--3959, 2020.

\bibitem{diehl2005realtime}
M. Diehl, H. G. Bock, and J. P. Schl\"oder, ``A real-time iteration scheme for nonlinear optimization in optimal feedback control,'' \emph{SIAM J. Control Optim.}, vol. 43, no. 5, pp. 1714--1736, 2005.

\bibitem{gros2020rti}
S. Gros, M. Zanon, R. Quirynen, A. Bemporad, and M. Diehl, ``From linear to nonlinear MPC: Bridging the gap via the real-time iteration,'' \emph{Int. J. Control}, vol. 93, no. 1, pp. 62--80, 2020.

\bibitem{zavala2009advancedstep}
V. M. Zavala and L. T. Biegler, ``The advanced-step NMPC controller: Optimality, stability and robustness,'' \emph{Automatica}, vol. 45, no. 1, pp. 86--93, 2009.

\bibitem{vaswani2017attention}
A. Vaswani \emph{et al.}, ``Attention is all you need,'' in \emph{Proc. NeurIPS}, 2017, pp. 5998--6008.

\bibitem{nie2023time}
Y. Nie, N. H. Nguyen, P. Sinthong, and J. Kalagnanam, ``A time series is worth 64 words: Long-term forecasting with Transformers,'' in \emph{Proc. ICLR}, 2023.

\end{thebibliography}

\end{document}
